% !PRINTED 454
\weberpage{454}

This theorem, which will render us very important services in what follows, is almost self-evident; for the greatest common divisor of two functions $F(x)$ and $f(x)$ can be found by rational computation, and is therefore also contained in $\Omega$. But since $f(x)$ has no divisor contained in $\Omega$ except itself or a constant (i.\,e. a quantity contained in $\Omega$), this greatest common divisor must therefore be either a constant or $f(x)$ itself.

It follows from this theorem that an irreducible function can never have multiple factors; for otherwise it would have to have a common divisor with its derivative $f'(x)$; thus $f'(x)$ would have to be divisible by $f(x)$, which is impossible, since the degree of $f'(x)$ is lower than that of $f(x)$.

We can also give Theorem I the following expression:

\gesperrt{II. If $f(x)$ is irreducible, and $F(x)$ vanishes for one root of $f(x) = 0$, then it also vanishes for all other roots of $f(x) = 0$.}

In particular we can conclude from this:

\gesperrt{III. If $F(x)$ is of lower degree than the irreducible function $f(x)$, and $F(x)$ vanishes for a root of $f(x)$, then $F(x)$ must vanish identically, i.\,e. all coefficients of $F(x)$ must be zero.}

Integral rational functions of several variables, $f(x, y, z \ldots)$, whose coefficients belong to the field $\Omega$, are likewise called functions in $\Omega$. These too are designated as \gesperrt{reducible} and \gesperrt{irreducible} functions, according as they are decomposable or not decomposable into several factors which are themselves functions in $\Omega$, and actually contain at least some of the variables $x, y, z \ldots$.

Here, however, attention must be called to a distinction.

Among the functions of more than one independent variable there are some that are not decomposable at all into factors of lower degree which depend rationally on the variables; and again others that are indeed decomposable, but only into factors whose coefficients are not contained in $\Omega$; and finally, thirdly, functions that are decomposable into several functions in $\Omega$; only the last kind shall we call reducible in $\Omega$, while those of the second kind are described as \gesperrt{decompos-}
%CONT>>
